\section{Reprodução dos resultados}
\label{sec:reproducao}

O repositório\novo{~\cite{mqttsecure2026repo}} contém o código, os dados brutos de cada campanha
(\texttt{results/}) e os programas de análise. Os passos para repetir as medições
num Raspberry~Pi estão em \texttt{docs/raspberry-pi.md}; o procedimento e o
significado de cada ficheiro, em \texttt{docs/evaluation.md}.

\begin{itemize}[leftmargin=*]
  \item \textbf{Refazer as tabelas e figuras a partir dos dados brutos:}
    \path{tools/analysis/analyze.py} \path{results/<campanha>} e, nesta pasta,
    \texttt{make dados PI=<campanha>} \novo{\texttt{[DESKTOP=<campanha>]}} seguido de \texttt{make}.
  \item \textbf{Repetir as medições:} \texttt{loadgen.py campaign} com a mesma
    semente e o mesmo \texttt{params.toml} (guardados em cada campanha).
\end{itemize}

A Tabela~\ref{tab:campanhas} identifica as campanhas de que vêm os resultados\novo{: a
do Raspberry~Pi, usada em todo o relatório, e a de referência, usada só na
Secção~\ref{sec:referencia}}.

\begin{table}[ht]
  \centering
  \caption{Campanhas usadas neste relatório.}
  \label{tab:campanhas}
  \footnotesize
  \begin{tabularx}{\textwidth}{@{}L{3.6cm}Y@{}}
    \toprule
    Nome & \texttt{pi5-20261113} (\texttt{results/pi5-20261113/}) \\
    Plataforma & Raspberry~Pi~5, 8~GB (\emph{broker} e \emph{gateway}); computador
      de secretária (dois agentes de sensores e um atacante) \\
    Data & 12--13 de novembro de 2026, 23:37 a 04:32 (UTC) \\
    \emph{Commit} & \texttt{79c4878}, árvore limpa \\
    Semente & 20261201 \\
    Execuções & 120 de custo (12 configurações $\times$ 10); subida de carga com
      3 repetições por configuração (427 patamares); matriz de ataques
      $3 \times 156$; 39 cenários; \emph{microbenchmarks} com 10 repetições em
      três variantes de aceleração \\
    Testes no mesmo \emph{commit} & \texttt{results/validation/pi5-20261113/} \\
    \midrule
    \novo{Nome} & \novo{\texttt{desktop-20261115} (\texttt{results/desktop-20261115/})} \\
    \novo{Plataforma} & \novo{Computador de secretária (AMD Ryzen~7 3700X): tudo na
      mesma máquina, com núcleos separados (Secção~\ref{sec:metodo})} \\
    \novo{Data} & \novo{Ataques, \emph{microbenchmarks} e cenários: 15--16 de novembro de
      2026, 23:40 a 01:10 (UTC); custo e subida de carga, repetidos: 16 de
      novembro, 06:23 a 10:20 (UTC)} \\
    \novo{\emph{Commit}} & \novo{\texttt{2f5f62c}, árvore limpa; o código de medição é
      igual ao de \texttt{79c4878} (só mudou \texttt{tools/analysis/analyze.py})} \\
    \novo{Semente} & \novo{20261201} \\
    \novo{Execuções} & \novo{As mesmas da campanha do Raspberry~Pi, com mais um patamar
      na subida de carga (544 patamares); \emph{microbenchmarks} com e sem AES-NI.
      A primeira tentativa da etapa de custo parou na execução 70 e fica em
      \texttt{7a-falhada/}, fora da análise} \\
    \bottomrule
  \end{tabularx}
\end{table}

\subsection*{Plataforma de referência}

\novo{Tabelas completas da campanha \texttt{desktop-20261115}
(Secção~\ref{sec:referencia}). A latência foi medida pela interface local e não é
comparável com a do Raspberry~Pi.}

\geradoRef{desktop/summary}
\geradoRef{desktop/saturation}
\geradoRef{desktop/criterion}
